\documentclass[12pt,a4paper]{article}
\usepackage{amsmath,amssymb,amsthm}
\usepackage{booktabs}
\usepackage{hyperref}
\usepackage{geometry}
\geometry{margin=2.5cm}
\usepackage{enumitem}

\newtheorem{theorem}{Theorem}[section]
\newtheorem{lemma}[theorem]{Lemma}
\newtheorem{corollary}[theorem]{Corollary}
\newtheorem{proposition}[theorem]{Proposition}
\theoremstyle{definition}
\newtheorem{definition}[theorem]{Definition}
\theoremstyle{remark}
\newtheorem{remark}[theorem]{Remark}

\title{The Weinberg Angle from Recognition Science:\\
$\sin^2\theta_W = (3 - \varphi)/6 \approx 0.230$}

\author{Jonathan Washburn\\
Recognition Science Research Institute\\
Austin, Texas\\
\texttt{washburn.jonathan@gmail.com}}

\date{March 2026}

\begin{document}
\maketitle

\begin{abstract}
We derive the weak mixing angle $\theta_W$ from the Recognition
Science (RS) $D$-cube structure.  The formula:
$\sin^2\theta_W = (D - \varphi)/(2D) = (3 - \varphi)/6
\approx 0.2303$, where $D = 3$ is the spatial dimension and
$\varphi = (1+\sqrt{5})/2$ is the golden ratio.  The
$\overline{\mathrm{MS}}$ value at $M_Z$ (PDG 2024) is
$\sin^2\hat\theta_W = 0.23121 \pm 0.00003$, giving agreement
to $0.4\%$.  The on-shell value $1 - m_W^2/m_Z^2 \approx 0.2231$
is 3.2\% from the RS prediction; the RS formula most naturally
corresponds to the $\overline{\mathrm{MS}}$ scheme at $M_Z$.
The RS formula connects electroweak mixing directly to the golden
ratio and the dimension of space: $\sin^2\theta_W$ is the
fraction of gauge space not filled by the $\varphi$-ladder
(structural derivation; open problem to establish from the RS
gauge lattice in full rigor).  We prove $0 < \sin^2\theta_W < 1/2$,
derive the $W/Z$ mass ratio $m_W/m_Z = \cos\theta_W$, and show
that only $D = 3$ reproduces the observed value.
\end{abstract}

%% ============================================================
\section{Introduction}
\label{sec:intro}
%% ============================================================

The Weinberg angle $\theta_W$~\cite{Weinberg1967,Glashow1961}
parameterizes the mixing between the $SU(2)_L$ and $U(1)_Y$ gauge
fields in the electroweak theory.  The physical photon and $Z$ boson
are:
\begin{align}
  A_\mu &= B_\mu \cos\theta_W + W_\mu^3 \sin\theta_W, \\
  Z_\mu &= -B_\mu \sin\theta_W + W_\mu^3 \cos\theta_W.
\end{align}
The mixing parameter
$\sin^2\theta_W = g'^2/(g_2^2 + g'^2)$ is a free parameter in
the Standard Model.  Recognition Science~\cite{RS_Axioms} derives it
from the dimension $D = 3$ and the golden ratio $\varphi$.

%% ============================================================
\section{Recognition Science Framework}
\label{sec:framework}
%% ============================================================

Recognition Science derives all physics from a single primitive,
the \emph{recognition cost functional} $J(x)$, uniquely determined
by three axioms.

\begin{definition}[RS Cost Functional]
For $x > 0$: $J(x) = \tfrac{1}{2}(x + x^{-1}) - 1$.
\end{definition}

\noindent\textbf{Axiom A1 (Normalization):} $J(1) = 0$.\\
\textbf{Axiom A2 (Recognition Composition Law, RCL):}
$J(xy) + J(x/y) = 2J(x)J(y) + 2J(x) + 2J(y)$ for all $x, y > 0$.\\
\textbf{Axiom A3 (Calibration):} $J''_{\log}(0) = 1$, where
$J_{\log}(t) := J(e^t)$.

\begin{theorem}[Cost Uniqueness]
\label{thm:unique}
The unique continuous solution to \textup{(A1)--(A3)} is
$J(x) = \tfrac{1}{2}(x + x^{-1}) - 1$.
\end{theorem}

\begin{proof}[Proof sketch]
Setting $f(t) = J_{\log}(t) + 1$ and rewriting \textup{(A2)} gives the
d'Alembert equation $f(s+t) + f(s-t) = 2f(s)f(t)$.  By the Acz\'el
classification theorem~\cite{Aczel1966}, the unique continuous solution
with $f(0) = 1$ and $f''(0) = 1$ is $f(t) = \cosh t$.
\end{proof}

\begin{theorem}[$\varphi$ Forcing]
\label{thm:phi}
The golden ratio $\varphi = (1+\sqrt{5})/2$ is the unique positive real
satisfying the self-similarity equation forced by the RCL\@.  All stable
masses lie on the $\varphi$-ladder: $m(r) = E_{\rm coh}\cdot\varphi^r$,
where $E_{\rm coh} = \varphi^{-5}$ in RS units and $r \in \mathbb{Z}$.
\end{theorem}

\noindent Spatial dimension $D = 3$ is forced by the gap-45 synchronization
$\operatorname{lcm}(8,45) = 360$, which uniquely requires $D = 3$ faces on
the fundamental cube.

%% ============================================================
\section{Derivation from the $D$-Cube}
\label{sec:derivation}
%% ============================================================

\begin{definition}[$D$-Cube Gauge Counting]
\label{def:gauge}
The $D$-cube has $2D$ faces.  In RS, these $2D$ face-pairs are
the gauge directions---each face corresponds to a generator of
the gauge group.  The total gauge space has dimension $2D$.
\end{definition}

\begin{proposition}[Hypercharge Fraction]
\label{prop:hyper}
Of the $2D$ gauge directions, the $\varphi$-ladder structure
``fills'' $D + \varphi$ of them (the $D$ spatial directions
plus $\varphi$ units of self-similar overlap).  The remaining
fraction:
\begin{equation}
  \frac{2D - (D + \varphi)}{2D} = \frac{D - \varphi}{2D}
\end{equation}
corresponds to the $U(1)_Y$ (hypercharge) component of the
gauge space.
\end{proposition}

\begin{remark}[Status of Proposition~\ref{prop:hyper}]
The claim that $D + \varphi$ directions are ``filled'' by the
$\varphi$-ladder is a structural hypothesis: the $D = 3$ spatial
directions fill $D$ gauge degrees of freedom, and the golden-ratio
self-similar scaling contributes an effective $\varphi$ units of
overlap between the $SU(2)_L$ and $U(1)_Y$ sectors.  A fully
rigorous derivation of this effective filling fraction from the
RS lattice gauge theory is an open problem.  The formula's agreement
with experiment ($0.4\%$) supports its structural basis.
\end{remark}

\begin{proposition}[Weinberg Angle --- RS Prediction]
\label{thm:weinberg}
\begin{equation}
  \boxed{\sin^2\theta_W = \frac{D - \varphi}{2D}
  = \frac{3 - \varphi}{6} \approx 0.2303.}
\end{equation}
\end{proposition}

\begin{proof}
The electroweak mixing parameter is the hypercharge fraction of
the total gauge space.  By Proposition~\ref{prop:hyper}:
\begin{equation}
  \sin^2\theta_W = \frac{D - \varphi}{2D}
  = \frac{3 - (1+\sqrt{5})/2}{6}
  = \frac{(5 - \sqrt{5})/2}{6}
  = \frac{5 - \sqrt{5}}{12}.
\end{equation}
Numerically:
\begin{equation}
  \sin^2\theta_W = \frac{5 - 2.236}{12}
  = \frac{2.764}{12} = 0.2303.
\end{equation}
\end{proof}

\begin{remark}
The formula has no free parameters: $D = 3$ and $\varphi$ are
both forced by the RS framework.  The denominator $2D = 6$ is
the number of cube faces; the numerator $D - \varphi$ is the
``unfilled'' gauge space.
\end{remark}

%% ============================================================
\section{Properties}
\label{sec:properties}
%% ============================================================

\begin{remark}[Bounds]
\label{thm:bounds}
From the formula $(3-\varphi)/6$:
\begin{enumerate}[label=(\roman*)]
  \item $0 < \sin^2\theta_W < 1/2$: since $1 < \varphi < 3$,
  $0 < (3-\varphi)/6 < 1/3 < 1/2$.
  \item $\sin^2\theta_W \neq 1/4$: $(3-\varphi)/6 = 1/4$ requires
  $\varphi = 3/2$, but $\varphi = (1+\sqrt{5})/2 \neq 3/2$.
  This rules out tree-level $SU(5)$ as a coincidence.
  \item $|\sin^2\theta_W - 0.230| = 0.0003 < 0.002$.
\end{enumerate}
\end{remark}

\begin{corollary}[W/Z Mass Ratio]
\label{thm:wz}
\begin{equation}
  \frac{m_W}{m_Z} = \cos\theta_W
  = \sqrt{1 - \frac{D - \varphi}{2D}}
  = \sqrt{\frac{D + \varphi}{2D}} \approx 0.877.
\end{equation}
\end{corollary}

\begin{proof}
$\cos^2\theta_W = 1 - \sin^2\theta_W = 1 - (D-\varphi)/(2D)
= (2D - D + \varphi)/(2D) = (D + \varphi)/(2D)$.
$(3 + 1.618)/6 = 4.618/6 = 0.7697$, so
$\cos\theta_W = \sqrt{0.7697} = 0.877$.
The experimental value is $m_W/m_Z = 80.369/91.188 = 0.8813$, a
$0.5\%$ deviation from the RS prediction.
\end{proof}

\begin{corollary}
$m_W < m_Z$ (since $\cos\theta_W < 1$).
\end{corollary}

%% ============================================================
\section{Running of $\sin^2\theta_W$}
\label{sec:running}
%% ============================================================

\begin{remark}[Scheme dependence and RS comparison]
\label{rem:scheme}
The RS formula $(3-\varphi)/6 = 0.2303$ gives a ``bare''
value at the recognition scale.  In quantum field theory,
$\sin^2\theta_W$ runs with energy and depends on the
renormalization scheme.  The three most commonly quoted values
are distinct:
\begin{center}
\renewcommand{\arraystretch}{1.2}
\begin{tabular}{lcc}
\toprule
\textbf{Scheme / Scale} & \textbf{Measured (PDG 2024)} &
\textbf{RS deviation} \\
\midrule
On-shell: $1 - m_W^2/m_Z^2$ & $\approx 0.2231$~\cite{PDG2024} & $3.2\%$ \\
$\overline{\mathrm{MS}}$, $M_Z$ & $0.23121 \pm 0.00003$~\cite{PDG2024} & $0.4\%$ \\
Effective leptonic, $M_Z$ & $0.23153 \pm 0.00016$~\cite{PDG2024} & $0.5\%$ \\
\bottomrule
\end{tabular}
\end{center}
The RS formula $0.23033$ agrees well with the
$\overline{\mathrm{MS}}$ scheme value at $M_Z$ (0.4\% deviation)
and the effective leptonic value (0.5\%).  The on-shell scheme
value ($\approx 0.2231$, from $m_W = 80.369~\mathrm{GeV}$,
$m_Z = 91.188~\mathrm{GeV}$) is 3.2\% from the RS prediction,
suggesting that the RS tree-level formula most naturally
corresponds to the $\overline{\mathrm{MS}}$ or effective
scheme at $M_Z$.  The identification of the RS formula's
scheme correspondence is a structural open problem.
\end{remark}

%% ============================================================
\section{Experimental Comparison}
\label{sec:comparison}
%% ============================================================

\begin{center}
\renewcommand{\arraystretch}{1.3}
\begin{tabular}{lccc}
\toprule
\textbf{Quantity} & \textbf{RS} &
\textbf{Exp.\ (PDG 2024)} & \textbf{Dev.} \\
\midrule
$\sin^2\hat\theta_W(M_Z)_{\overline{\rm MS}}$ & $0.23034$ & $0.23122\pm 0.00003$ & $0.4\%$ \\
$\sin^2\theta_W$ (on-shell) & $0.23034$ & $\approx 0.2231$ & $3.2\%$ \\
$m_W/m_Z$ & $0.877$ & $0.881$ & $0.5\%$ \\
$\sin^2\theta_W$ (eff.\ leptonic) & $0.23034$ & $0.23153\pm 0.00016$ & $0.5\%$ \\
\bottomrule
\end{tabular}
\end{center}

%% ============================================================
\section{Dimension Dependence}
\label{sec:dimension}
%% ============================================================

\begin{proposition}[Only $D = 3$ Works]
\label{prop:dimension}
The formula $(D - \varphi)/(2D)$ gives:
\begin{center}
\renewcommand{\arraystretch}{1.2}
\begin{tabular}{ccc}
\toprule
$D$ & $\sin^2\theta_W$ & Status \\
\midrule
2 & $(2 - \varphi)/4 = 0.0955$ & Ruled out ($\gg\!10\sigma$) \\
3 & $(3 - \varphi)/6 = 0.2303$ & Matches ($\overline{\rm MS}$: $0.4\%$) \\
4 & $(4 - \varphi)/8 = 0.2977$ & Ruled out ($\gg\!10\sigma$) \\
\bottomrule
\end{tabular}
\end{center}
Only $D = 3$ is compatible with the observed value.
\end{proposition}

%% ============================================================
\section{Comparison with Other Predictions}
\label{sec:other}
%% ============================================================

\begin{center}
\renewcommand{\arraystretch}{1.3}
\begin{tabular}{lccc}
\toprule
\textbf{Framework} & \textbf{$\sin^2\theta_W$} &
\textbf{Free params} & \textbf{Error ($\overline{\rm MS}$)} \\
\midrule
$SU(5)$ GUT (tree, $M_\mathrm{GUT}$) & $3/8 = 0.375$ & 0 & $62\%$ \\
$SO(10)$ GUT (run to $M_Z$) & $\sim 0.21$ & 0 & $\sim 9\%$ \\
Standard Model (fitted) & $0.23122$ & 1 & $0$ (by fit) \\
RS (zero parameters) & $(3-\varphi)/6 = 0.23034$ & 0 & $0.4\%$ \\
\bottomrule
\end{tabular}
\end{center}

\begin{remark}
GUT theories predict $\sin^2\theta_W = 3/8$ at the unification
scale, which must run down to $\approx 0.23$ at $M_Z$.  RS
predicts $0.2303$ directly at the electroweak scale, without
invoking unification or running from a high scale.
\end{remark}

%% ============================================================
\section{Predictions and Falsifiers}
\label{sec:predictions}
%% ============================================================

\begin{enumerate}
  \item \textbf{$\sin^2\theta_W = (3-\varphi)/6 = 0.2303$}:
  Zero-parameter prediction at the electroweak scale.

  \item \textbf{$m_W/m_Z = \sqrt{(3+\varphi)/6} = 0.877$}:
  Derived from the same formula.

  \item \textbf{Falsifier}: A precision measurement showing
  the running of $\sin^2\theta_W$ is inconsistent with
  $(3-\varphi)/6$ as the tree-level value would challenge the
  formula.

  \item \textbf{No GUT unification}: RS predicts no high-energy
  scale at which the gauge couplings unify in the $SU(5)$ sense.
\end{enumerate}

%% ============================================================
\section{Conclusion}
\label{sec:conclusion}
%% ============================================================

The Weinberg angle is derived from the $D = 3$ cube geometry:
$\sin^2\theta_W = (3-\varphi)/6 = (5-\sqrt{5})/12 \approx 0.23034$,
matching the experimental $\overline{\rm MS}$ value at $M_Z$
($0.23122 \pm 0.00003$) to $0.4\%$ with zero free parameters.
The formula connects electroweak mixing directly to the golden ratio
and the dimension of space.  The on-shell scheme value ($\approx
0.2231$) shows a $3.2\%$ discrepancy, suggesting the RS tree-level
formula corresponds most naturally to the $\overline{\rm MS}$ scheme
at $M_Z$.  A rigorous derivation of the hypercharge-fraction formula
from the RS lattice gauge theory remains an open problem.

%% ============================================================
\begin{thebibliography}{99}

\bibitem{Aczel1966}
J.~Acz\'el, \emph{Lectures on Functional Equations and Their Applications},
Academic Press, New York (1966).

\bibitem{RS_Axioms}
J.~Washburn, ``The Algebra of Reality: A Recognition Science Derivation
of Physical Law,'' \emph{Axioms} \textbf{15}(2), 90 (2026).
\texttt{doi:10.3390/axioms15020090}

\bibitem{Weinberg1967}
S.~Weinberg, ``A Model of Leptons,'' Phys.\ Rev.\ Lett.\
\textbf{19}, 1264 (1967).

\bibitem{PDG2024}
S.~Navas \textit{et al.}\ (PDG), ``Review of Particle Physics,''
Phys.\ Rev.\ D \textbf{110}, 030001 (2024).

\bibitem{Glashow1961}
S.~L.~Glashow, ``Partial-Symmetries of Weak Interactions,''
Nucl.\ Phys.\ \textbf{22}, 579 (1961).

\end{thebibliography}

\end{document}
